% appendices/B_dsc_proofs.tex
% Faithful port of the "Extended Proofs" part (Sec. S2) of the ECCV 2026
% supplementary: ../sources/eccv2026_domain_feature_collapse/supplementary.tex.
% Geometric proofs for Domain-Sensitivity Collapse (DSC): the linear-model
% illustration, Theorem (distance failure under variance--discriminability
% mismatch), and Proposition (MSP/Energy insensitivity). The theorem and
% proposition STATEMENTS live in the chapter body (Chapter on DSC,
% \Cref{thm:dsc:dist-fail} / \Cref{prop:dsc:msp-insens}); this appendix supplies
% the proofs. v1/Appendix B (mutual-information / information-bottleneck proofs)
% does NOT carry over.
%
% NOTE on the number gate (make lint): lint.py strips math line-by-line and does
% not recognize equation* or multi-line environments, so every display equation
% here is kept on a single physical line inside a stripped environment.
\chapter{Geometric Proofs for Domain-Sensitivity Collapse}
\label{app:dsc-proofs}

This appendix supplies the proofs of the two geometric results stated in the
chapter on Domain-Sensitivity Collapse (DSC): the distance-failure theorem
(\Cref{thm:dsc:dist-fail}) and the MSP/Energy insensitivity proposition
(\Cref{prop:dsc:msp-insens}). We first record a linear-model illustration that
makes concrete, in the simplest possible setting, why single-domain supervised
training relegates the discriminative signal to a low-dimensional subspace and
suppresses the directions along which out-of-distribution (OOD) inputs differ
from in-distribution (ID) inputs. We then prove an orthogonal-energy lemma that
both proofs rely on, and finally the theorem and proposition themselves.

Throughout, we use the notation established in the DSC chapter. We write
$\mathcal{S}_{\mathrm{cls}} = \mathrm{span}\{v_1,\ldots,v_k\}$ for the
class-discriminative subspace spanned by the top eigenvectors of the ID feature
covariance, $P$ for the orthogonal projector onto $\mathcal{S}_{\mathrm{cls}}$,
and $P_\perp = I - P$ for the projector onto its orthogonal complement. The
defining DSC condition is that ID features carry negligible energy outside the
dominant subspace,
\begin{equation} \mathbb{E}\bigl[\|P_\perp z\|_2^2\bigr] \le \tau^2 \quad \text{for } z = f_\theta(x_{\mathrm{ID}}), \label{eq:dsc:A1} \end{equation}
and we assume the analogous bound with constant $\tau_{\mathrm{OOD}}$ holds for
OOD features; under the dominant-subspace--matching hypothesis the two tail
constants are comparable, and we write $\tau$ for either when no confusion
arises.

\section{Linear-Model Illustration}
\label{sec:dsc-proofs:linear}

For completeness we record the linear objective underlying the informal
mechanism sketched in the chapter, which shows in closed form why a regularized
supervised model discards the directions that distinguish ID from OOD inputs
under a domain shift.

\begin{lemma}[Linear-model collapse]\label{lem:dsc:linear}
Let $y \in \{\pm 1\}$ and let the input decompose as $x = y\,a + s$, where
$a$ is a fixed class direction and $s \perp y$ is a nuisance component carrying
no label information. Consider the linear representation $z = Ax$ scored by a
linear head $w$, trained under the regularized logistic objective
\begin{equation} \min_{A,\,w}\;\mathbb{E}\bigl[\log(1 + e^{-y\,w^\top\! A x})\bigr] + \tfrac{\lambda}{2}\bigl(\|w\|_2^2 + \|A\|_F^2\bigr). \label{eq:dsc:toy-loss} \end{equation}
Then there is an optimum $A^\star$ whose row space lies in $\mathrm{span}(a)$.
Consequently, for any OOD shift $\delta \in \mathrm{span}(a)^\perp$, the
induced representation shift vanishes: $A^\star \delta = 0$, so
$z_{\mathrm{OOD}} - z_{\mathrm{ID}} = A^\star \delta = 0$.
\end{lemma}

\begin{proof}
The logistic loss in~\eqref{eq:dsc:toy-loss} depends on $A$ only through the
scalar margin $w^\top\! A x$, and hence only through the row vector $w^\top\! A$
acting on $x$. Decompose any row of $A$ into its component in $\mathrm{span}(a)$
and its component in $\mathrm{span}(a)^\perp$. Because $s \perp y$, the data
distribution of $x = y a + s$ couples the label to $x$ only through the
$\mathrm{span}(a)$ component; the orthogonal component of $w^\top\! A$ therefore
contributes nothing to the expected margin in either direction (its
contribution averages to a label-independent term that the logistic loss is
minimized by setting to zero in expectation). Thus the data-fit term is
invariant under projecting every row of $A$ onto $\mathrm{span}(a)$, whereas the
Frobenius regularizer $\tfrac{\lambda}{2}\|A\|_F^2$ is strictly decreased by
removing any nonzero orthogonal component. Hence at least one minimizer
$A^\star$ has all rows in $\mathrm{span}(a)$, i.e.\ row space contained in
$\mathrm{span}(a)$.

For such an $A^\star$ and any $\delta \in \mathrm{span}(a)^\perp$ we have
$A^\star \delta = 0$, because each row of $A^\star$ is orthogonal to $\delta$.
An OOD perturbation confined to $\mathrm{span}(a)^\perp$ therefore produces no
change in the representation, $z_{\mathrm{OOD}} - z_{\mathrm{ID}} = 0$, and no
score built on $z$ can separate the two.
\end{proof}

\begin{remark}
\Cref{lem:dsc:linear} is the linear shadow of the general phenomenon: the
$\ell_2$ regularizer (or, in deep networks, the implicit bias of gradient
descent toward simpler solutions together with weight decay) strips away every
representation direction that does not reduce the supervised loss. When the
domain-sensitive directions along which OOD inputs differ are label-irrelevant
in the single-domain training distribution---precisely the case
$\delta \in \mathrm{span}(a)^\perp$---those directions are suppressed, and the
representation becomes blind to the shift. In the multi-dimensional setting this
manifests as the concentration of variance in $\mathcal{S}_{\mathrm{cls}}$ and
the tail-energy bound~\eqref{eq:dsc:A1}.
\end{remark}

\section{Orthogonal-Energy Lemma}
\label{sec:dsc-proofs:ortho}

Both proofs reduce ambient distances to in-subspace distances at the cost of a
controlled tail term. The following lemma quantifies that cost.

\begin{lemma}[Orthogonal-energy bound]\label{lem:dsc:ortho-energy}
For any feature vectors $z, z'$,
\begin{equation} \bigl|\,\|z - z'\|_2^2 - \|Pz - Pz'\|_2^2\,\bigr| = \|P_\perp(z - z')\|_2^2 \le 2\|P_\perp z\|_2^2 + 2\|P_\perp z'\|_2^2 . \label{eq:dsc:ortho-energy} \end{equation}
In particular, if both $z$ and $z'$ satisfy the tail-energy
condition~\eqref{eq:dsc:A1} in expectation with constant $\tau$, then
\begin{equation} \mathbb{E}\Bigl[\,\bigl|\,\|z - z'\|_2^2 - \|Pz - Pz'\|_2^2\,\bigr|\,\Bigr] \;\le\; 4\tau^2 . \label{eq:dsc:ortho-energy-exp} \end{equation}
\end{lemma}

\begin{proof}
Because $P$ and $P_\perp$ are complementary orthogonal projectors, the
Pythagorean identity gives, for the difference vector $u = z - z'$, the
decomposition $\|u\|_2^2 = \|Pu\|_2^2 + \|P_\perp u\|_2^2$, where
$Pu = Pz - Pz'$ by linearity of $P$. Subtracting $\|Pu\|_2^2$ from both sides
and taking absolute values yields the equality in~\eqref{eq:dsc:ortho-energy}.
For the inequality, linearity of $P_\perp$ and the parallelogram bound
$\|p - q\|_2^2 \le 2\|p\|_2^2 + 2\|q\|_2^2$ give
$\|P_\perp(z - z')\|_2^2 = \|P_\perp z - P_\perp z'\|_2^2 \le 2\|P_\perp z\|_2^2 + 2\|P_\perp z'\|_2^2$.
Taking expectations and applying~\eqref{eq:dsc:A1} to each term (with constant
$\tau$ for both ID and OOD features) bounds the right-hand side by
$2\tau^2 + 2\tau^2 = 4\tau^2$, which is~\eqref{eq:dsc:ortho-energy-exp}.
\end{proof}

\section{Proof of the Distance-Failure Theorem}
\label{sec:dsc-proofs:thm}

\begin{proof}[Proof of \Cref{thm:dsc:dist-fail}]
Write the $k$-nearest-neighbour score for a test point $z$ as
$S_k(z) = \|z - z_{(k)}\|_2$, where $z_{(k)}$ is the $k$-th closest training
point in Euclidean distance. For a fixed training set, $S_k$ is $1$-Lipschitz in
$z$: for any $z, z'$, the triangle inequality gives
$|S_k(z) - S_k(z')| \le \|z - z'\|_2$, since the same training point $z_{(k)}$
can be used as a feasible (not necessarily optimal) neighbour for both
arguments. Hence the Lipschitz constant of the statistic satisfies
$L_k \le 1$.

\emph{Step~$1$ (projection to the discriminative subspace).}
By \Cref{lem:dsc:ortho-energy}, replacing the ambient distance $\|z - z'\|_2$ by
the in-subspace distance $\|Pz - Pz'\|_2$ perturbs the squared score by at most
$\|P_\perp(z - z')\|_2^2$, whose expectation is bounded by $4\tau^2$ under
\eqref{eq:dsc:A1} for both ID and OOD inputs. Passing from squared to unsquared
distances and using $1$-Lipschitzness, the corresponding change in the $L_1$
statistic of $S_k$ is of the same order, controlled by the tail constant
$\tau$.

\emph{Step~$2$ (matching inside the class-discriminative subspace).}
By hypothesis the ID-vs-OOD separation concentrates in the low-variance
directions $\{v_j : j \in \mathcal{J}\}$ with
$\lambda_j/\lambda_1 \le \rho \ll 1$. Consequently the residual ID-vs-OOD
discrepancy that remains \emph{inside} the high-variance part of
$\mathcal{S}_{\mathrm{cls}}$ is small; write
$\varepsilon = \mathbb{E}\bigl[\|P z_{\mathrm{ID}} - P z_{\mathrm{OOD}}\|_2\bigr]$
for this residual. Because the separating directions are down-weighted by the
anisotropic ID geometry, the projected ID and OOD score distributions are within
$\varepsilon$ of one another in $W_1$.

\emph{Step~$3$ (Wasserstein bound).}
Combining the two steps, the distributions of $S_k(z_{\mathrm{ID}})$ and
$S_k(z_{\mathrm{OOD}})$ differ by at most $L_k\,\varepsilon$ from the in-subspace
matching of Step~$2$, plus the null-space perturbation of Step~$1$, which
contributes $4 L_k \tau^2$ via the expected orthogonal-energy
bound~\eqref{eq:dsc:ortho-energy-exp}. Since the $W_1$ distance is bounded by the
expected absolute difference of the (coupled) score values, we obtain
\begin{equation} W_1\!\bigl(S_{\mathrm{KNN}}(z_{\mathrm{ID}}),\; S_{\mathrm{KNN}}(z_{\mathrm{OOD}})\bigr) \;\le\; L_k\,\varepsilon + 4 L_k\,\tau^2 , \label{eq:dsc:knn-fail} \end{equation}
with $L_k \le 1$, which is the claimed bound.
\end{proof}

\paragraph{On the constants.}
The Lipschitz constant $L_k \le 1$ is exactly the $1$-Lipschitzness of the
$k$-NN distance established at the start of the proof. The coefficient $4$
multiplying $\tau^2$ is inherited verbatim from the expected orthogonal-energy
bound~\eqref{eq:dsc:ortho-energy-exp} of \Cref{lem:dsc:ortho-energy}. The bound
is informative precisely in the DSC regime: when $\varepsilon \approx 0$ (OOD
data projects onto $\mathcal{S}_{\mathrm{cls}}$ almost identically to ID) and
$\tau^2 \ll 1$ (strong tail suppression), the right-hand side
of~\eqref{eq:dsc:knn-fail} collapses to $\approx 4\tau^2$, so the ID and OOD
score distributions are nearly indistinguishable and distance-based detection
fails. Conversely, when the features are isotropic the tail constant grows
($\tau \approx \sqrt{d}$) and the bound becomes vacuous, exactly the regime in
which distance scores should and do succeed.

\paragraph{Mahalanobis distance also fails.}
The same argument disposes of the shrinkage-regularized Mahalanobis score
$M_c(z) = (z - \mu_c)^\top (\Sigma_c + \lambda I)^{-1}(z - \mu_c)$, which is
sometimes believed to recover the suppressed directions by up-weighting
low-variance coordinates. Decomposing $z - \mu_c$ into its $P$ and $P_\perp$
parts, the contribution from $\mathcal{S}_{\mathrm{cls}}^\perp$ is bounded above
by $(1/\lambda)\|P_\perp(z - \mu_c)\|_2^2$, which is $O(\tau^2/\lambda)$ under
\eqref{eq:dsc:A1}. The shrinkage term $\lambda I$ caps the amplification of the
null space, so when~\eqref{eq:dsc:A1} holds for both ID and OOD inputs the
Mahalanobis score---and any post-hoc distance score operating on the same
low-rank features, including ZCA whitening, which is algebraically equivalent to
Mahalanobis with the global covariance estimate---cannot leverage the
discriminative tail. This is a statement about the whole class of distance-based
scores acting on a collapsed representation, not a claim that no scoring rule
whatsoever could recover the signal.

\section{Proof of the MSP/Energy Insensitivity Proposition}
\label{sec:dsc-proofs:prop}

We first record that the logit-based scores are Lipschitz in the logits. Let the
logits be $\ell(x) = W z(x) + b$, and recall the energy score
$S_{\mathrm{Energy}}(\ell) = -T \log \sum_c \exp(\ell_c / T)$, taken with unit
temperature $T = 1$.

\begin{proof}[Proof of \Cref{prop:dsc:msp-insens}]
\emph{Lipschitzness of the score.}
The gradient of the energy score is
$\nabla_\ell S_{\mathrm{Energy}} = -\,\mathrm{softmax}(\ell)$, so
$\|\nabla_\ell S_{\mathrm{Energy}}\|_2 = \|\mathrm{softmax}(\ell)\|_2 \le \|\mathrm{softmax}(\ell)\|_1 = 1$,
and thus $S_{\mathrm{Energy}}$ is $1$-Lipschitz in $\ell$; that is,
$L_S \le 1$ for the energy score. For the maximum-softmax-probability score, the
softmax Jacobian $\mathrm{diag}(p) - p p^\top$ has operator norm at most $2$,
giving $L_S \le 2$ for MSP.

\emph{Decomposition of the logit shift.}
The logit difference between an OOD and an ID input is
$\ell(x_{\mathrm{OOD}}) - \ell(x_{\mathrm{ID}}) = W\bigl(z_{\mathrm{OOD}} - z_{\mathrm{ID}}\bigr)$,
and decomposing the feature shift as
$z_{\mathrm{OOD}} - z_{\mathrm{ID}} = P(z_{\mathrm{OOD}} - z_{\mathrm{ID}}) + P_\perp(z_{\mathrm{OOD}} - z_{\mathrm{ID}})$
splits it into an in-subspace part $W P(z_{\mathrm{OOD}} - z_{\mathrm{ID}})$ and
a null-space part $W P_\perp(z_{\mathrm{OOD}} - z_{\mathrm{ID}})$.

\emph{Bounding the two parts.}
By the dominant-subspace--matching hypothesis, the in-subspace shift is small:
$\|P(z_{\mathrm{OOD}} - z_{\mathrm{ID}})\|_2 \le \varepsilon$, so the first term
satisfies
$\|W P(z_{\mathrm{OOD}} - z_{\mathrm{ID}})\|_2 \le \|W\|_{\mathrm{op}}\,\varepsilon$.
By head insensitivity $\|W P_\perp\|_{\mathrm{op}} \le \eta$ together with the
tail-energy condition~\eqref{eq:dsc:A1} applied to both ID and OOD features, the
second term satisfies
$\|W P_\perp(z_{\mathrm{OOD}} - z_{\mathrm{ID}})\|_2 \le \eta\,(\|P_\perp z_{\mathrm{OOD}}\|_2 + \|P_\perp z_{\mathrm{ID}}\|_2) \le \eta\,(\tau_{\mathrm{OOD}} + \tau) = O(\eta\tau)$
in expectation.

\emph{Composition.}
Composing the score's Lipschitz bound with the logit-shift bound and using that
the $W_1$ distance is dominated by the expected absolute score difference under
a coupling, we obtain
\begin{equation} W_1\!\bigl(S(\ell(x_{\mathrm{ID}})),\; S(\ell(x_{\mathrm{OOD}}))\bigr) \;\le\; \|W\|_{\mathrm{op}}\,\varepsilon + L_S\,\eta\tau , \label{eq:dsc:logit-fail} \end{equation}
for $S \in \{S_{\mathrm{Energy}}, S_{\mathrm{MSP}}\}$, with $L_S \le 1$ for the
energy score and $L_S \le 2$ for MSP. This is the claimed bound.
\end{proof}

\paragraph{On the constants.}
Every constant in~\eqref{eq:dsc:logit-fail} is directly measurable from the
trained model. The factor $\|W\|_{\mathrm{op}}$ is the largest singular value of
the classifier head, which is $O(1)$ under standard weight decay; the Lipschitz
constants $L_S$ are intrinsic to the score function ($L_S \le 1$ for energy,
$L_S \le 2$ for MSP); and the head insensitivity
$\eta = \|W P_\perp\|_{\mathrm{op}}$ is the spectral norm of $W$ restricted to
$\mathcal{S}_{\mathrm{cls}}^\perp$, which approaches zero as neural-collapse
alignment concentrates the row space of $W$ in $\mathcal{S}_{\mathrm{cls}}$. As
with the distance-failure theorem, the bound is informative exactly in the DSC
regime $\varepsilon \ll 1$, $\tau \ll 1$, and becomes vacuous when the features
are isotropic or the OOD inputs differ strongly inside the class subspace---the
regimes in which logit-based scores should succeed.
